% status: ZERO
% Chapter 36 of the second edition. Plan: docs/STRUCTURE.md. Rules: AUTHORING.md.
% Measurements: research/measurements/2026-09-24-m1-part7.md (replay, selftest);
% Hermon's own records at commit 2a3fd52 (the developers' measurements).
\chapter{Streaming Experts From Disk: A Measured Case Study}\label{ch:expert-streaming}

\begin{ChapterIntent}
Hermon's expert pager streams the experts of a 30-billion-parameter
mixture-of-experts model from an NVMe drive into a fixed cache. This chapter
follows its measured development: hit rate turned out to be the whole lever, an
attractive queue-depth optimization delivered nothing, the bottleneck moved
from I/O to compute and back as SIMD and threads landed, and parallelizing
experts exposed a result that depended on the thread count. It is a case study
in predicting, measuring and being wrong --- and in saying exactly what is
integrated and what is not.
\end{ChapterIntent}

\OpeningSection{Tokens per second from an SSD}

The first time Hermon's expert pager ran against a real checkpoint, it
reported 10.08 tokens per second, reading experts from a consumer NVMe drive at
8,331~MiB/s. No single consumer drive sustains that rate. The container holding
the experts of Qwen3-30B-A3B was 16.35~GiB and the machine had 30~GB of RAM, so
the operating system's page cache had served most of the reads: the benchmark
was timing memory and reporting storage (Hermon commit \texttt{be0288a},
2026-08-25; the developers' measurement). Its direct-I/O switch could not have
worked. \texttt{O\_DIRECT} constrains the alignment of the destination buffer as
well as the file offset and length, and the cache's blocks were aligned to 64
bytes. With page-aligned blocks, and the file switched to direct I/O only after
its unaligned header had been read, the same benchmark reported 2.12 tokens per
second with no cache hits and 7.50 with 71\% of expert reads hitting the cache
--- ceilings set by the drive alone, because the benchmark moved the experts'
bytes and computed nothing with them.

Those two corrections set the pattern for the twelve days of development this
chapter follows. Every step made a prediction, measured it, and more than once
found the prediction or the measurement wrong. The pager is also a small
enough system to check from the outside: this chapter's replay regenerates the
benchmark's routing trace from its seed and reproduces every hit rate the
developers published, to the printed digit (derived by simulation; record
\texttt{2026-09-24-m1-part7}).

\NapkinSection{Napkin math: hit rate, bytes per miss, tokens per second}

An MoE layer routes each token to $k$ of its experts. If an expert's weights
occupy $r_\ell$ bytes in layer $\ell$ and a fraction $h$ of expert reads hit the
cache, a token reads $(1-h)\,k\sum_\ell r_\ell$ bytes from the drive. Its expert
arithmetic is $F = 2\cdot 3\,d\,d_{\mathrm{ff}}\,k\,L$ FLOPs for $L$ MoE layers
of three $d \times d_{\mathrm{ff}}$ matrices (gate, up, down). If reads and
arithmetic do not overlap, the time per token for the expert layers is
\begin{equation}
t_{\mathrm{tok}} \approx \frac{(1-h)\,k\sum_{\ell} r_\ell}{\beta_{\mathrm{io}}}
 + \frac{F}{\pi_{\mathrm{eff}}},
\label{eq:expert-stream-time}
\end{equation}
with $\beta_{\mathrm{io}}$ the drive's sustained read bandwidth and
$\pi_{\mathrm{eff}}$ the FLOP rate the expert kernels achieve; with perfect
prefetching the sum becomes a maximum.

Qwen3-30B-A3B has $L = 48$, $k = 8$ of 128 experts, $d = 2{,}048$ and an expert
width $d_{\mathrm{ff}} = 768$ (its published configuration). In llama.cpp's
Q4\_K\_M format the gate and up projections are Q4\_K and the down projection is
Q6\_K in 24 layers and Q4\_K in the other 24, so an expert occupies 3,059,712 or
2,654,208 bytes (2.92 or 2.53~MiB). A token that misses everything reads
$8 \times 24 \times (3{,}059{,}712 + 2{,}654{,}208) = 1.10$~GB (1.02~GiB), and
its expert arithmetic is 3.62~GFLOP (derived). At 2.3~GB/s the reads take
0.48~s --- 2.1 tokens per second with no hits and 7.3 at a 71\% hit rate ---
while the arithmetic takes 1.6~s at the 2.2~GFLOP/s of a scalar loop and
0.095~s at the 38~GFLOP/s eight threads later reached. Which term dominates
depends on the code, not on the model, and the rest of this chapter is the two
terms trading places.

\section{The problem: more experts than memory}

A mixture-of-experts model computes like its active parameters and must be
stored like its total parameters (\Chref{mixture-of-experts}). Qwen3-30B-A3B
activates 3.3 of its 30.5 billion parameters per token; its experts alone come
to 16.35~GiB in 4-bit form, more than a 16~GB laptop can hold beside everything
else. At the frontier the gap is two orders of magnitude: Kimi~K3 has 2.8
trillion parameters, 16 of its 896 experts active per token, and ships its
experts in MXFP4 \cite{moonshot2026kimik3} --- well over a terabyte, beyond
the memory of any laptop or desktop.

Streaming is the answer two independent C implementations converged on, and
Hermon's design followed them. WARP (formerly WASTE) keeps a 27.28~GB trunk
resident, streams Kimi~K3's experts from a 982~GB container with one aligned
read per expert into a bounded LRU cache, and reports 0.6 tokens per second on
a 64~GB M5~Pro laptop \cite{sqliteai2026warp}. Kimi-K3-in-C runs the same model
in 8.24~GB of RAM and reports that, at a fixed memory budget, keeping the dense
trunk resident beat giving the memory to an expert cache by 1.69 times
\cite{khan2026kimik3c}. Both are their authors' measurements. WARP adds a
warning that shaped Hermon's cache: an expert cache larger than physical RAM was
eight times slower, because its hits had become page faults.

\section{An expert store}

Hermon stores the experts in a container file with one aligned record per
(layer, expert): gate, up and down concatenated, so that a miss is exactly one
\texttt{pread} at an offset computed from the layer and expert index
(\texttt{crates/hermon-kernels/csrc/experts.c} at \texttt{2a3fd52}). The layout is
a clustered index in the database sense --- the same bytes as the GGUF file,
physically reordered for one access pattern at the cost of a second copy. The
real checkpoint broke the first version of it. Q4\_K\_M mixes formats, so records
differ between layers; padding every record to the largest would have cost
7.1\% of the container and 15\% more bytes per read in half the layers
(3.06 against 2.65~MB, derived), in a design whose premise is that read
bandwidth binds. The second version stores a stride per layer.

The cache is not a new allocator. It is the same block pool Hermon's paged KV
cache uses --- blocks with reference counts and a free list --- with an LRU list
over resident experts. Two bugs in it are worth knowing because no status check
would have caught either. A resident block must hold a reference owned by the
store itself: without it, a block whose caller had finished returned to the
free list while the residency map still pointed at it, and a later request for
that expert silently received another expert's weights with a success code.
And acquisition must be all-or-nothing: a partial failure that left some blocks
pinned leaked them, until the cache could no longer allocate. The rule the
reference count now encodes is short --- one reference means cached and
evictable, more than one means in use --- and the blocks one layer is computing
with are pinned against eviction by the next miss.

\section{Hit rate is the lever}

The developers replayed a 64-token routing trace against caches holding 2 to
50\% of the 6,144 expert records, with two routing distributions generated by
the benchmark: uniform, and a power law over a fixed random ordering of each
layer's experts \cite[section~K6.1]{hermon2026kerneldesign}:

\begin{center}
\small
\begin{tabular}{lrrrr}
\toprule
Routing & Cache (GiB) & Hit rate & Throughput (as published) & Tokens/s ceiling \\
\midrule
uniform & 0.35 & 0.0\% & 2,373 MiB/s & 2.12 \\
uniform & 1.75 & 6.3\% & 2,422 MiB/s & 2.31 \\
uniform & 4.38 & 22.0\% & 2,413 MiB/s & 2.76 \\
uniform & 8.75 & 44.0\% & 2,411 MiB/s & 3.84 \\
skewed & 1.75 & 20.4\% & 2,416 MiB/s & 2.71 \\
skewed & 4.38 & 49.9\% & 2,370 MiB/s & 4.22 \\
skewed & 8.75 & 71.3\% & 2,409 MiB/s & 7.50 \\
\bottomrule
\end{tabular}
\end{center}

The flat column is the finding. The drive delivered the same rate whatever the
cache size or routing, so the tokens-per-second ceiling was a function of the
hit rate alone --- Equation~\ref{eq:expert-stream-time} with the compute term
absent. The replay reproduces the hit-rate column exactly: its routing trace
is regenerated from the benchmark's own generator and seed, and its cache
follows the pager's pinning and eviction rules. It also recounts the bytes. The
published reads were the number of misses times the largest record, which
overstates them by 7.0 to 7.1\% because half the records are smaller; the
developers found and fixed that accounting five days later
\cite{hermon2026batchio}. Recounted, the drive delivered about 2.3~GB/s, not
2.4 --- a correction that leaves every ceiling in the table unchanged, because
those were timed per token rather than derived from bytes.

Two cautions keep the table honest. The skewed trace is synthetic: a Zipf law
of exponent one over a per-layer permutation, not a recording of Qwen3's
router. Trained routers are pushed toward balanced long-run load
(\Chref{mixture-of-experts}), but consecutive tokens do reuse experts:
Eliseev and Mazur built an LRU expert cache for Mixtral-8x7B on that
observation and ran it at 2.3 tokens per second on an RTX~3060 with 2-bit
experts \cite{eliseev2023moeoffload}. Qwen3-30B-A3B's own router is not
recorded here; a smaller model's is, at the end of this section. And a 64-token
trace that starts from an empty cache understates the steady state, because up
to a quarter of its reads are first touches.

The replay can run the same generators for 2,048 tokens and score only the
second half, and there it found something the published table could not show
(Figure~\ref{fig:ch36-hit-rate}). Warm, LRU hit 47.6\% of reads with uniform
routing and 77.2\% with skewed routing at the 8.75~GiB cache, against 44.0 and
71.3 cold. But at the two smallest caches, 2 and 5\% of the experts, LRU hit
nothing even with skewed routing, while simply pinning the experts that the
first half of the trace had used most hit 18.0 and 32.9\% of the second half's
reads (derived by simulation). The access pattern is a loop: every token visits
the 48 layers in order, so between two visits to a layer the other 47 read 376
distinct records, and an LRU cache smaller than that has evicted the layer's
experts before the next token comes back to them --- the textbook failure of LRU
on cyclic scans. Frequency survives the loop. At
8.75~GiB with skewed routing the learned frequency policy reached 84.2\% and an
offline optimum (Belady's, which knows the future) 91.3\%; at 2.3~GB/s those
are ceilings of 9.2 tokens per second for LRU, 13 for the frequency policy and
24 for the optimum (derived). LRU was the right first policy --- simple, and
correct under pinning --- and the replay says where the next gains are.

\EngineFigure{ch36-hit-rate}{Expert-cache hit rate against cache size for the
pager benchmark's traces (Qwen3-30B-A3B shape: 6,144 expert records, 8 of 128
per layer). Cold LRU over 64 tokens is what the developers measured, which the
replay reproduces exactly; the warm curves score the second half of a
2,048-token replay. LRU hits nothing until the cache holds more than the 376
records the other layers read between two visits to a layer, while with skewed
routing a frequency policy learned from the first half keeps most of the
optimum (Belady's) at every size. Both routings are the
benchmark's synthetic distributions, not a recorded router; derived by
simulation.}{fig:ch36-hit-rate}

A recorded router gives the same verdict. IBM's Granite 3.1 3B-A800M routes each
token to 8 of 40 experts in each of its 32 layers, small enough to run on the
M1, and llama.cpp recorded its choices over 2,705 tokens generated greedily from
eight varied prompts --- code, arithmetic, a story, history, translation and two
passages of technical prose (measured; record \texttt{2026-09-26-m1-part8}). The
router was far from uniform: a layer's 8 most used experts took 43\% of its
selections, where uniform routing would give 20\%, and consecutive tokens shared
46\% of their experts in a layer, against 20\%. Replayed as one session through
caches of 5 to 50\% of the 1,280 expert records, scoring the second half, the
recorded trace reproduced the loop: at 5 and 10\%, below the 248 records the
other 31 layers read between two visits to a layer, LRU hit nothing, while the
most used experts of the first half hit 16 and 27\% and Belady's optimum 22 and
42. At a quarter of the records LRU recovered to 46\%, and at half to 77\%, level
with the frequency policy (derived by simulation from the measured trace). Real
routing has the locality the synthetic trace assumed; it does not rescue LRU from
the layer cycle.

\section{Queue depth: a falsified hypothesis}

With the pager bound by the drive, the next lever looked obvious. Every miss
was a synchronous read at queue depth one; the design took a consumer NVMe
drive to deliver 5 to 7~GB/s with 32 requests in flight against about 2.4 with
one; and the eight misses of a layer are known at the same instant. It
predicted two to three times the throughput from submitting them together. The
pager was restructured to classify a layer's hits and misses first and submit
the misses as one batch through \texttt{io\_uring}, with an environment variable
to turn it off as the control. Measured against the sequential path in the same
run, the mean over ten configurations was 1,108~MiB/s sequential and
1,145~MiB/s batched --- indistinguishable \cite{hermon2026batchio}.

The prediction had applied a specification to the wrong request size.
Queue-depth figures describe small random reads, which are latency-bound: each
waits for the device, so overlapping them fills idle time. A block-size sweep
on the same drive located the change: 128~KiB reads reached 494~MB/s, 1~MiB
reads 1.2~GB/s and 3~MiB reads 1.3~GB/s. A 2.5 to 2.9~MiB expert record already
streams near the device's limit, leaving nothing for a second request to fill.
The absolute numbers are lower than K6.1's because the machine was running
another inference server and its disk was 98\% busy during these runs --- the
comparison is matched, the absolute rates are not clean. The batching code was
kept --- it costs nothing measurable and pays the moment records get smaller ---
with its off switch as the standing control; the performance claim was
withdrawn. The restructuring also produced two bugs, both caught
by tests: a block claimed for a miss but not yet installed was orphaned when a
later read failed, and a router that named the same expert twice in one layer
claimed two blocks for one slot. The replay shows the benchmark's own trace had
contained such repeats --- 4 in the uniform trace and 9 in the skewed one, from
a de-duplication loop that skipped its first entry after a collision --- which
the single-pass cache had absorbed as hits (derived); the same commit fixed the
loop.

\section{Computing on streamed experts}

Until early September the pager moved bytes nobody used: its matrix--vector
kernel knew F32, F16, BF16 and Q8\_0, and the checkpoint's experts were Q4\_K and
Q6\_K. The developers implemented both super-block formats from the
specification rather than borrowing llama.cpp's code, which made the bit
packing the whole risk, and checked it three ways that catch different errors:
a test that chooses logical values first and packs them, so that packer and
kernel cannot share a mistake; an independent NumPy dequantizer run against
real bytes from the container, agreeing to the last bit; and the distribution
of the real weights (mean near zero, standard deviation about 0.023), which a
misread layout would distort by orders of magnitude \cite{hermon2026kquant}.

The composite an MoE layer needs per routed expert,
$\mathbf{y} \mathrel{+}= w\,\mathbf{W}_{\mathrm{down}}(\mathrm{silu}(\mathbf{W}_{\mathrm{gate}}\mathbf{x}) \odot \mathbf{W}_{\mathrm{up}}\mathbf{x})$,
then ran straight out of cached records. Its first result that mattered was
not a speed: over the same trace and weights, the output checksum was
identical at every cache size and under both routing distributions. Eviction
changes when bytes arrive, never what is computed --- the property that makes a
bounded cache safe to put on a request path.

\section{The bottleneck moves}

With real arithmetic the benchmark ran at 0.58 to 0.62 tokens per second
against the pager's 2.1: the scalar K-quant loop reached about 2.2~GFLOP/s, and
compute now cost three to four times the I/O \cite{hermon2026expertffn}. The
obvious remedy was SIMD, which had given the attention and KV-write paths 8
and 19 times. The first vectorized version was slower than scalar --- 0.62
against 0.72 tokens per second --- because it expanded quantized values into
16-element vectors and ended every short dot product with a horizontal
reduction that cost more than the multiply-adds it summed. Folding each
sub-block's scale into the expansion, so that a 256-element super-block is
reduced once, gave 0.83 and 0.73 tokens per second against a scalar control of
0.59 and 0.59 on the developers' Zen~3 desktop: about 1.3 times end to end and
1.7 times on the arithmetic once the I/O in the trace is set aside
\cite{hermon2026simd}. The developers recorded why it was not 8 times:
their kernel still converts every quantized value to float, whereas llama.cpp
multiplies in the integer domain and applies scales at the end.

These runs used a two-token trace per configuration, so they are single
measurements of a short workload, not rates to plan with. What they establish
is the order of the terms: for a while, arithmetic, not storage, set the pace.

\section{Parallel experts and a determinism bug}

A layer's eight experts are independent, so the next step ran them across
threads spawned by the benchmark (the library itself spawns none). With the
cache warm, so that the measurement is arithmetic and not I/O, one expert
layer took 8.43, 7.06, 3.82 and 1.97~ms on 1, 2, 4 and 8 threads --- 8.96 to
38.25~GFLOP/s, 4.27 times at eight threads --- sublinear because each expert
streams 2.9~MiB of weights through the caches on every call
\cite{hermon2026parallel}.

The first version was also wrong in a way no correctness test would see. Each
worker kept a private accumulator and the accumulators were summed in worker
order; with experts assigned by stride, worker order is not expert order, so
the order of the floating-point additions depended on the thread count. One
thread gave a checksum of $-11.581605$ and four gave $-11.581564$: both correct to
tolerance, neither reproducible from the other, and Hermon's release gate is
bit-for-bit equality. The fix was the one Hermon's attention combine step
already used --- one slot per expert, summed by expert index --- after which the
checksum was identical at 1, 2, 4 and 8 threads. Only a checksum compared
across thread counts found the bug; \Chref{fast-wrong-answers} makes that kind
of check a discipline.

At eight threads the expert arithmetic of a token takes about 0.095~s, while its
reads take 0.48~s with no hits and 0.14~s at a 71\% hit rate (derived with
Equation~\ref{eq:expert-stream-time} from 1.10~GB per token at 2.3~GB/s; the
developers' own summary divided GiB by GB/s and put these at 0.43 and 0.125~s).
The arc closed where it began: storage-bound, then compute-bound, then
storage-bound again, with the hit rate deciding everything and the next gain
in reading fewer bytes rather than computing them faster.

\section{What is and is not integrated}

At commit \texttt{2a3fd52} the pager is a library (LIBRARY): the store, the
cache, K-quant expert kernels and a safe Rust interface whose pin sets release
their experts on every return path, exercised by the kernel crate's own tests and
benchmark and by no request path. Hermon's roadmap says so directly: expert
streaming is not yet an inference feature, and runtime routing and an
end-to-end comparison of generated tokens remain
\cite[phase K4--K7]{hermon2026roadmap}. The figures above are ceilings for the
expert layers --- attention, the router and the rest of the trunk are not in
them --- and the design's two headline targets are unmet (DESIGN): at least
0.5 tokens per second end to end for a Kimi-K3-class checkpoint on a 64~GB
machine, and a 20\% gain from letting one allocator trade KV-cache capacity
against expert capacity (\Chref{unified-memory}). The library's self-test did
reproduce on different hardware: on the Apple M1 used for this book's
measurements it passed all 121,165 checks, plain and under the undefined-behaviour
sanitizer, including the expert-cache and K-quant cases (measured; record
\texttt{2026-09-24-m1-part7}).

\LedgerSection

Streaming the experts of a mixture-of-experts model from an SSD:

\begin{ConstraintLedger}
Memory capacity & buys & A model whose experts exceed RAM runs at all; the cache holds only the experts in use. \\
Memory bandwidth & spends & Every miss reads a 2.5--2.9~MiB record from a drive delivering about 2.3~GB/s, more than twenty times below the host's DRAM. \\
Latency & spends & Tenths of a second per token of reads at realistic hit rates, before any arithmetic. \\
Correctness & risks & A residency or eviction bug returns another expert's weights silently; summation order can depend on thread count. \\
\end{ConstraintLedger}

\LabSection{replay a routing trace against an expert cache}

\LabIntent{In \texttt{code/labs/ch36-expert-streaming/}, a replay of the pager
benchmark: it regenerates the routing trace from the benchmark's generator and
seed, runs it through a pinned LRU cache, and must reproduce the published hit
rates before anything else is trusted (the book's replay does, in
\texttt{research/measurements/2026-09-24-m1-part7/}). Then it compares LRU with
an offline optimum (Belady) and a static cache of the most-used experts, from a
cold and a warm start; measures the reader's own SSD with caching disabled; and
turns hit rates into tokens-per-second ceilings with
Equation~\ref{eq:expert-stream-time}. Break it by removing the pin on the
current layer's experts and watch a miss evict an expert still in use.}

\HappenedSection{the whole chapter}

The case is the chapter; what generalizes is its record of being wrong. Four
predictions were tested. The pager was predicted to be bound by the drive, and
was. Queue depth was predicted to give two to three times, and gave nothing.
SIMD was predicted to give a large multiple, and gave 1.3 times until the
reason was understood. Threads gave 4.27 times at eight, less than linear, for
a stated reason. Five numbers were wrong before they were right: a throughput
that exceeded what any single drive can deliver (the page cache), a byte count
7\% high (misses times the largest record), a relative-error tolerance that
failed a correct kernel because one output had cancelled to near zero, a
checksum that depended on the thread count, and a final summary that mixed
GiB with GB. Each was caught by a check that compared the number with something
independent of the code that produced it --- the device's physical limit, a
recount, a scale-aware error measure, the same run at another thread count, a
re-derivation. That is the discipline the rest of the book asks for, visible in
one small system over twelve days.

\FrontierSection{2026-09-24}

Streaming has moved from research systems to hobbyist-scale engines that run
trillion-parameter models on laptops at fractions of a token per second
\cite{sqliteai2026warp,khan2026kimik3c}. The alternatives trade the same bytes
differently: keep experts in host memory and compute them on the CPU
(\Chref{offloading}), or predict the next layer's experts early enough to
overlap their reads with computation --- by applying the next layer's gate to
the current hidden state \cite{eliseev2023moeoffload}, or by training the
model to route one layer ahead \cite{hwang2023pregated}. Each prediction scheme
changes only timing if the real router still decides. The open measurement is
the one this chapter could not make: how much locality real routers have on
real text, which decides every hit rate above.

\ExercisesSection

\begin{enumerate}
\item \emph{Derivation.} Rewrite Equation~\ref{eq:expert-stream-time} for a
  pager that prefetches layer $\ell+1$'s experts during layer $\ell$'s
  arithmetic with prediction recall $\rho$, and find the recall at which the
  eight-thread configuration becomes compute-bound at a 71\% hit rate.
\item \emph{Prediction.} A drive sustains 7~GB/s for 3~MiB reads. Predict the
  hit rate Qwen3-30B-A3B's experts need for 10 tokens per second of expert
  layers at eight threads, with and without overlap.
\item \emph{Implementation.} Extend the lab with a cache that pins each layer's
  $n$ most-used experts and runs LRU over the rest; find the split that
  maximizes the warm hit rate on the skewed trace.
\item \emph{Falsification.} The replay's LRU hit nothing with caches of 122
  and 307 records. Show that an LRU cache of fewer than $(L-1)\,k$ records (376
  here) scores zero on any trace that visits the layers in a fixed cycle, then
  construct a trace --- bursty routing, where an expert stays hot
  for a few tokens and then goes cold --- on which LRU beats the learned
  frequency policy. What property of the trace decides between them?
\end{enumerate}
